\section{Compiler output and lowering details}

Compilation returns:
\begin{lstlisting}[language={}]
CompileResult {
  gates: CircuitGate[]
  qubitCount: number
  logicalQubitCount: number
  parsed: ParsedCommand[]
  log: string[]
  tokenMap: Record<string, number>
  processParams: ProcessParam[]
  returnValues: ReturnValue[]
}
\end{lstlisting}

Each accepted instruction entry records \code{op}, original \code{raw} text,
\code{inputs}, \code{outputs}, positional \code{args}, optional normalized
radian \code{phase}, \code{reverse}, and \code{noParameterSubstitution}.
Each emitted gate records \code{id}, \code{type}, ordered \code{step},
\code{targets}, \code{controls}, optional \code{phase}, and source text.

After child expansion, unused symbolic holes are compacted to dense qubit
indices. Internal RESET gates count for simulation but are filtered from the
visible circuit. UI logical width prefers RETURNVALS count, then exposed PARAMS
count, then physical simulator width.

